\documentclass[11pt]{article}
\usepackage[a4paper,margin=25mm]{geometry}
\usepackage{amsmath,amssymb,amsthm}
\usepackage{booktabs}
\usepackage{caption}
\captionsetup[table]{position=top,skip=6pt}
\usepackage[hidelinks]{hyperref}
\hypersetup{pdftitle={H34\_16: y\textasciicircum 3-y\textasciicircum 2+y=x\textasciicircum 4+x+2 has no integer solutions}}
\newtheorem{theorem}{Theorem}
\newtheorem{lemma}[theorem]{Lemma}
\newtheorem{proposition}[theorem]{Proposition}
\newcommand{\Z}{\mathbb{Z}}
\newcommand{\Q}{\mathbb{Q}}
\newcommand{\F}{\mathbb{F}}
\title{\textbf{H34\_16}\\[4pt] The equation $y^3-y^2+y=x^4+x+2$ has no integer solutions}
\author{}
\date{}
\begin{document}
\maketitle
\vspace{-8ex}
\begin{center}
Size $h=34$, length $l=12$, identifier E1154 (with $y$ replaced by $-y$)
\end{center}

\begin{theorem}\label{thm:main}
The equation
\begin{equation}\label{eq:main}
y^3-y^2+y=x^4+x+2
\end{equation}
has no integer solutions.
\end{theorem}

The proof is an obstruction coming from the cubic Hilbert reciprocity law over the field $K=\Q(\omega)$, where $\omega^2+\omega+1=0$. To a hypothetical integer solution we attach the values of $20$ auxiliary polynomials $H_0,\dots,H_{19}$ with coefficients in $\Z[\omega]$, and a sum $B$ of cubic Hilbert symbols of these values, which is recorded by a weighted graph. By the Hilbert reciprocity law, the local contributions $B_p$ of all primes $p$ add up to $0$ modulo $3$. We show that $B_p=0$ for $p\neq2,3$ (Section~\ref{sec:good}), and that $B_2=2$ and $B_3=0$ (Section~\ref{sec:local}). This is a contradiction.

Two standard facts from class field theory are quoted without proof: the explicit formula for the cubic Hilbert symbol at the places not above $3$, and the Hilbert reciprocity law. Everything else is proved below, up to explicit finite computations: polynomial identities and resultants, which can be checked by expanding both sides or by polynomial division, and the local computations of Section~\ref{sec:local}, which are finite enumerations of residue classes.

\section{\texorpdfstring{Cubic Hilbert symbols over $\Q(\omega)$}{Cubic Hilbert symbols over Q(w)}}\label{sec:symbols}

\paragraph{The field.} Let $K=\Q(\omega)$. Its ring of integers is $\Z[\omega]$, and every element of $K$ can be written uniquely as $A+B\omega$ with $A,B\in\Q$. We write
\[
\langle A,B\rangle=A+B\omega,
\]
also when $A$ and $B$ are polynomials in $x$, $y$ or $t$ with rational coefficients. Then $\langle A,B\rangle\langle C,D\rangle=\langle AC-BD,\ AD+BC-BD\rangle$, and the norm is $N(A+B\omega)=A^2-AB+B^2$. The units of $\Z[\omega]$ are $\pm\omega^k$.

\paragraph{Places.} The finite places of $K$ correspond to the non-zero prime ideals of $\Z[\omega]$. A prime $p\equiv1\pmod3$ has two places above it, both with residue field $\F_p$, and a prime $p\equiv2\pmod3$ is inert, with residue field $\F_{p^2}$. The prime $3$ is ramified: $\pi=1-\omega$ generates the unique prime ideal above $3$, and $\pi^2=-3\omega$, so $9=\omega\pi^4$. The only infinite place of $K$ is complex. For a finite place $v$ we denote by $K_v$ the completion of $K$ at $v$ and by $v(a)$ the normalized valuation of $a\in K_v^\times$, and we call $a$ a \emph{unit} at $v$ if $v(a)=0$. For a rational prime $p$, we write $v\mid p$ if $v$ is a place above $p$.

\paragraph{Hilbert symbols.} For a place $v$ of $K$ and $a,b\in K_v^\times$, the cubic Hilbert symbol $(a,b)_v$ is a cube root of unity, see \cite[Chapter~V, \S3]{N}. We write
\[
(a,b)_v=\omega^{\beta_v(a,b)},\qquad\beta_v(a,b)\in\Z/3\Z.
\]
We use the following three facts \cite[Chapter~V, \S3 and Chapter~VI, \S8]{N}.
\begin{itemize}
\item[(H1)] $\beta_v(a,bc)=\beta_v(a,b)+\beta_v(a,c)$, $\beta_v(a,b)=-\beta_v(b,a)$, and $\beta_v(a,c^3)=0$. In particular, $\beta_v(a,b)$ depends only on the classes of $a$ and $b$ in $K_v^\times/K_v^{\times3}$.
\item[(H2)] Let $v$ be a finite place not above $3$, let $q$ be the number of elements of its residue field, and let $\varpi$ be a uniformizer at $v$. Write $a=\varpi^eu$ and $b=\varpi^fw$ with units $u$ and $w$, and define $\chi(u)\in\Z/3\Z$ by $\bar u^{(q-1)/3}=\bar\omega^{\chi(u)}$, where the bar denotes the reduction to the residue field. Then
\[
\beta_v(a,b)=e\,\chi(w)-f\,\chi(u).
\]
In particular, $\beta_v(a,b)=0$ if $a$ and $b$ are both units.
\item[(H3)] (Hilbert reciprocity.) For $a,b\in K^\times$, $\beta_v(a,b)=0$ for all but finitely many places $v$, and $\sum_v\beta_v(a,b)=0$, where the sum is over all places of $K$. At the complex place, $\beta_v(a,b)=0$.
\end{itemize}
The general formula in (H2) contains an additional term $ef\chi(-1)$, which vanishes because $-1=(-1)^3$ is a cube. Formula (H2) fixes the normalization of the symbols; with the opposite sign convention all values $\beta_v$ change sign, and the proof below is the same.

\paragraph{The place above $2$.} The prime $2$ is inert, the residue field is $\F_4=\{0,1,\bar\omega,\bar\omega^2\}$, and $(q-1)/3=1$. Hence, for a unit $u=A+B\omega$, the value $\chi(u)$ is $0$, $1$ or $2$ according as $(A,B)$ is congruent to $(1,0)$, $(0,1)$ or $(1,1)$ modulo $2$. By (H2) and (H1), the class of $a=2^eu$ modulo cubes is described by the pair $(e\bmod3,\chi(u))\in(\Z/3\Z)^2$, and
\begin{equation}\label{eq:beta2}
\beta_v\bigl(a,b\bigr)=e\,\chi(w)-f\,\chi(u)\qquad\text{for }a=2^eu,\ b=2^fw.
\end{equation}
If $a$ is a rational integer, then its unit part is a rational odd number, and $\chi(u)=0$.

\paragraph{The place above $3$.}
\begin{lemma}\label{lem:M3}
Let $v$ be the place above $3$. Every non-zero element of $K_v$ is equal to $\pi^a\omega^b2^c(1+3\omega)^d$ times a cube, for unique $a,b,c,d\in\Z/3\Z$. If $\mathbf a$ and $\mathbf b$ are the exponent vectors $(a,b,c,d)$ of two elements $A$ and $B$, then $\beta_v(A,B)=\mathbf a^{\mathsf T}M_3\mathbf b$, where
\[
M_3=\begin{pmatrix}0&0&2&0\\0&0&1&2\\1&2&0&0\\0&1&0&0\end{pmatrix}\pmod3.
\]
\end{lemma}

\begin{proof}
For $\alpha\in\Z[\omega]$, the polynomial $g(s)=s+3s^2+3s^3-\alpha$ satisfies $g'(s)\equiv1\pmod\pi$, so by Hensel's lemma it has a root $s$ in the completion of $\Z[\omega]$, and then $(1+3s)^3=1+9(s+3s^2+3s^3)=1+9\alpha$. Hence every unit congruent to $1$ modulo $9$ is a cube in $K_v$, and the class of a unit modulo cubes depends only on its residue modulo $9$. There are $54$ units in $\Z[\omega]/(9)$, and a direct enumeration shows that their cubes are $\pm1$ modulo $9$, and that the $27$ products $\omega^b2^c(1+3\omega)^d$ with $b,c,d\in\{0,1,2\}$ are pairwise distinct modulo $9$, also after multiplication by $-1=(-1)^3$. Hence they represent all $27$ classes of units modulo cubes. Since $\pi^3$ is a cube, the exponent of $\pi$ is the valuation modulo $3$, and the first statement follows.

By (H1), it remains to compute $\beta_v$ on the $16$ ordered pairs of the elements $\pi$, $\omega$, $2$ and $1+3\omega$. Their norms are $3$, $1$, $4$ and $7$, hence, for each such pair, all the symbols at the places not above $2$, $3$ and $7$ vanish by (H2), and, by (H3), $\beta_v$ at the place above $3$ is minus the sum of the values at the place above $2$ and at the two places above $7$. These places correspond to $\bar\omega=2$ and $\bar\omega=4$ in $\F_7$. The classes $(e,\chi)$ of the four elements at these places are
\[
\begin{array}{c|ccc}
 & 2 & \bar\omega=2 & \bar\omega=4\\\hline
\pi & (0,2) & (0,0) & (0,2)\\
\omega & (0,1) & (0,2) & (0,2)\\
2 & (1,0) & (0,2) & (0,1)\\
1+3\omega & (0,2) & (1,0) & (0,0)
\end{array}
\]
where, at the places above $7$, $\chi(u)$ is defined by $\bar u^2=\bar\omega^{\chi(u)}$ in $\F_7$, and $7$ is a uniformizer. Formula (H2) and the sum over these three places give $M_3$.
\end{proof}

\section{The graph}\label{sec:graph}

Put $f(x,y)=y^3-y^2+y-x^4-x-2$, so that \eqref{eq:main} is the equation $f(x,y)=0$. The certificate consists of $20$ auxiliary polynomials $H_i$ and constants $c_i$, listed in Table~\ref{tab:vertices}, and of weights $e_{ij}\in\{1,2\}$ for $58$ pairs $i<j$, the \emph{edges}. Put $e_{ij}=0$ for the other pairs $i<j$. We also use $\rho_{ij}=e_{ij}$ for $i<j$ and $\rho_{ij}=-e_{ji}$ for $i>j$, computed modulo $3$; these numbers are listed in Table~\ref{tab:edges}.

\begin{table}[htbp]
\centering\small
\caption{The auxiliary polynomials $H_i$ and the constants $c_i$. Here $\langle A,B\rangle=A+B\omega$.}\label{tab:vertices}
\begin{tabular}{@{}rll@{}}
			\toprule
			$i$ & $H_i$ & $c_i$\\\midrule
		0 & $\langle y + 1,1\rangle$ & $\langle 4,0\rangle$\\
		1 & $\langle y - 1,-1\rangle$ & $\langle -4,-4\rangle$\\
		2 & $\langle x,0\rangle$ & $\langle 0,4\rangle$\\
		3 & $\langle x + 1,1\rangle$ & $\langle 0,2\rangle$\\
		4 & $\langle x + 1,0\rangle$ & $\langle 4,0\rangle$\\
		5 & $\langle x,1\rangle$ & $\langle 0,2\rangle$\\
		6 & $\langle x,-1\rangle$ & $\langle 2,0\rangle$\\
		7 & $\langle x - 1,-1\rangle$ & $\langle 1,0\rangle$\\
		8 & $\langle x + y,1\rangle$ & $\langle 4,0\rangle$\\
		9 & $\langle x,y - 2\rangle$ & $\langle 2,4\rangle$\\
		10 & $\langle - x + y - 2,- x^{2} - x - 1\rangle$ & $\langle 2,0\rangle$\\
		11 & $\langle y - 1,0\rangle$ & $\langle 0,1\rangle$\\
		12 & $\langle x - y + 1,0\rangle$ & $\langle 4,8\rangle$\\
		13 & $\langle x + y,0\rangle$ & $\langle 0,1\rangle$\\
		14 & $\langle y - 2,-1\rangle$ & $\langle 0,2\rangle$\\
		15 & $\langle y - 1,- x\rangle$ & $\langle 6,6\rangle$\\
		16 & $\langle - x + y,1 - x\rangle$ & $\langle 0,-6\rangle$\\
		17 & $\langle y^{2} - y + 1,- x y\rangle$ & $\langle 1,0\rangle$\\
		18 & $\langle x y + x + y^{2},x + y - 1\rangle$ & $\langle -1,-1\rangle$\\
		19 & $\langle - x^{2} + y^{2} + 1,- x^{2} - x y\rangle$ & $\langle -1,-1\rangle$\\
			\bottomrule
		\end{tabular}
\end{table}

\begin{table}[htbp]
\centering\small
\caption{The numbers $\rho_{ij}$. A dot denotes zero; $\rho_{ij}=e_{ij}$ for $i<j$, and $\rho_{ji}=-\rho_{ij}\bmod3$.}\label{tab:edges}
\setlength{\arraycolsep}{3.6pt}
$\begin{array}{r|*{20}{c}}
			\rho_{ij} & 0 & 1 & 2 & 3 & 4 & 5 & 6 & 7 & 8 & 9 & 10 & 11 & 12 & 13 & 14 & 15 & 16 & 17 & 18 & 19\\\hline
			0 & \cdot & \cdot & \cdot & \cdot & \cdot & \cdot & \cdot & \cdot & 1 & \cdot & \cdot & \cdot & \cdot & 1 & \cdot & \cdot & 1 & 2 & \cdot & \cdot\\
			1 & \cdot & \cdot & 2 & \cdot & 1 & 2 & 1 & \cdot & \cdot & \cdot & 2 & \cdot & 1 & 2 & \cdot & 2 & \cdot & 2 & \cdot & \cdot\\
			2 & \cdot & 1 & \cdot & \cdot & \cdot & \cdot & \cdot & \cdot & \cdot & 2 & \cdot & \cdot & \cdot & \cdot & 2 & \cdot & \cdot & \cdot & 2 & 2\\
			3 & \cdot & \cdot & \cdot & \cdot & \cdot & \cdot & \cdot & \cdot & \cdot & 2 & 2 & 1 & 1 & \cdot & \cdot & \cdot & 1 & \cdot & 1 & \cdot\\
			4 & \cdot & 2 & \cdot & \cdot & \cdot & \cdot & \cdot & \cdot & \cdot & \cdot & 2 & \cdot & 1 & \cdot & 2 & \cdot & \cdot & \cdot & \cdot & 2\\
			5 & \cdot & 1 & \cdot & \cdot & \cdot & \cdot & \cdot & \cdot & 2 & \cdot & \cdot & 2 & 2 & \cdot & \cdot & \cdot & \cdot & \cdot & 2 & 2\\
			6 & \cdot & 2 & \cdot & \cdot & \cdot & \cdot & \cdot & \cdot & 2 & 2 & 1 & \cdot & \cdot & \cdot & \cdot & \cdot & \cdot & \cdot & 2 & 2\\
			7 & \cdot & \cdot & \cdot & \cdot & \cdot & \cdot & \cdot & \cdot & \cdot & \cdot & \cdot & 1 & \cdot & \cdot & \cdot & \cdot & \cdot & 1 & \cdot & \cdot\\
			8 & 2 & \cdot & \cdot & \cdot & \cdot & 1 & 1 & \cdot & \cdot & 1 & \cdot & \cdot & \cdot & \cdot & \cdot & \cdot & 2 & 2 & \cdot & \cdot\\
			9 & \cdot & \cdot & 1 & 1 & \cdot & \cdot & 1 & \cdot & 2 & \cdot & \cdot & \cdot & 2 & \cdot & 2 & 2 & \cdot & \cdot & \cdot & \cdot\\
			10 & \cdot & 1 & \cdot & 1 & 1 & \cdot & 2 & \cdot & \cdot & \cdot & \cdot & 1 & \cdot & 2 & \cdot & \cdot & \cdot & 2 & \cdot & \cdot\\
			11 & \cdot & \cdot & \cdot & 2 & \cdot & 1 & \cdot & 2 & \cdot & \cdot & 2 & \cdot & \cdot & \cdot & \cdot & \cdot & 2 & \cdot & \cdot & \cdot\\
			12 & \cdot & 2 & \cdot & 2 & 2 & 1 & \cdot & \cdot & \cdot & 1 & \cdot & \cdot & \cdot & \cdot & \cdot & 2 & \cdot & \cdot & \cdot & \cdot\\
			13 & 2 & 1 & \cdot & \cdot & \cdot & \cdot & \cdot & \cdot & \cdot & \cdot & 1 & \cdot & \cdot & \cdot & \cdot & 1 & \cdot & \cdot & 1 & 1\\
			14 & \cdot & \cdot & 1 & \cdot & 1 & \cdot & \cdot & \cdot & \cdot & 1 & \cdot & \cdot & \cdot & \cdot & \cdot & 1 & \cdot & \cdot & \cdot & \cdot\\
			15 & \cdot & 1 & \cdot & \cdot & \cdot & \cdot & \cdot & \cdot & \cdot & 1 & \cdot & \cdot & 1 & 2 & 2 & \cdot & \cdot & 1 & 2 & \cdot\\
			16 & 2 & \cdot & \cdot & 2 & \cdot & \cdot & \cdot & \cdot & 1 & \cdot & \cdot & 1 & \cdot & \cdot & \cdot & \cdot & \cdot & \cdot & \cdot & 1\\
			17 & 1 & 1 & \cdot & \cdot & \cdot & \cdot & \cdot & 2 & 1 & \cdot & 1 & \cdot & \cdot & \cdot & \cdot & 2 & \cdot & \cdot & \cdot & \cdot\\
			18 & \cdot & \cdot & 1 & 2 & \cdot & 1 & 1 & \cdot & \cdot & \cdot & \cdot & \cdot & \cdot & 2 & \cdot & 1 & \cdot & \cdot & \cdot & 2\\
			19 & \cdot & \cdot & 1 & \cdot & 1 & 1 & 1 & \cdot & \cdot & \cdot & \cdot & \cdot & \cdot & 2 & \cdot & \cdot & 2 & \cdot & 1 & \cdot\\
		\end{array}$
\end{table}

Suppose that $(x,y)$ is an integer solution of \eqref{eq:main}, and write $H_i$ also for the value $H_i(x,y)\in\Z[\omega]$. We will show in Section~\ref{sec:nonzero} that all these values are non-zero. For every place $v$ of $K$ put
\begin{equation}\label{eq:B}
B_v=\sum_{i<j}e_{ij}\,\beta_v(H_i,H_j)+\sum_{i=0}^{19}\beta_v(H_i,c_i)\in\Z/3\Z,
\end{equation}
and, for a rational prime $p$, put $B_p=\sum_{v\mid p}B_v$.

\begin{proposition}\label{prop:reciprocity}
If all $H_i$ are non-zero, then $B_v=0$ for all but finitely many places $v$, and $\sum_pB_p=0$.
\end{proposition}

\begin{proof}
Apply (H3) to every pair $(H_i,H_j)$ with $e_{ij}\neq0$, multiply by $e_{ij}$, and apply it also to every pair $(H_i,c_i)$. Adding the results gives $\sum_vB_v=0$, where the complex place contributes $0$. Grouping the finite places above the same rational prime gives the claim.
\end{proof}

\section{The auxiliary values are non-zero}\label{sec:nonzero}

Let $(x,y)$ be an integer solution of \eqref{eq:main}. Since $1$ and $\omega$ are linearly independent over $\Q$, a value $H_i=\langle A,B\rangle$ vanishes if and only if $A(x,y)=B(x,y)=0$.
\begin{itemize}
\item For $i=0,1,3,5,6,7,8,14$, the coefficient $B$ of $H_i$ is a non-zero constant, and for $i=10$ it is $-(x^2+x+1)\neq0$. Hence these $H_i$ do not vanish.
\item If $H_9$, $H_{15}$ or $H_{16}$ vanished, then the two coefficients would give $(x,y)=(0,2)$, $(0,1)$ or $(1,1)$, respectively. But $f(0,2)=4$, $f(0,1)=-1$ and $f(1,1)=-3$.
\item The polynomials $H_2=x$, $H_4=x+1$, $H_{11}=y-1$, $H_{12}=x-y+1$ and $H_{13}=x+y$ have rational coefficients. Substituting $H_i=0$ into $f$ gives the following polynomials in the remaining coordinate $t$ (namely $y$ for $i=2,4,12$ and $x$ for $i=11,13$):
\[
\begin{array}{c|l}
i& f(x,y)\text{ on }H_i=0\\\hline
2,4&t^3-t^2+t-2\\
11&-t^4-t-1\\
12&-t^4+5t^3-7t^2+4t-2\\
13&-t^4-t^3-t^2-2t-2
\end{array}
\]
An integer root of such a polynomial divides its constant term, so the only candidates are $\pm1$ for $i=11$ and $\pm1,\pm2$ in the other rows, and substitution shows that none of them is a root.
\item The coefficient $A$ of $H_{17}$ is $y^2-y+1>0$. If $H_{18}=0$, then $B=x+y-1=0$, and then $A=xy+x+y^2=1$. If $H_{19}=0$, then $B=-x(x+y)=0$, and then $A=-x^2+y^2+1$ is equal to $y^2+1$ or to $1$.
\end{itemize}
Hence no $H_i$ vanishes at an integer solution.

\section{\texorpdfstring{The primes $p\neq2,3$}{The primes p other than 2 and 3}}\label{sec:good}

\paragraph{Charts.} We parametrize the curve $H_i=0$ by a variable $t$. For $i=0,1,8,10,11,13,14,15,16$, the polynomial $H_i$ has the form $y-r_i(x)$, and we put $\sigma_i(t)=(t,r_i(t))$. For $i=2,3,4,5,6,7,9,12$, it has the form $x-r_i(y)$, and we put $\sigma_i(t)=(r_i(t),t)$. For example, $\sigma_0(t)=(t,-1-\omega)$ and $\sigma_9(t)=(-(t-2)\omega,t)$. The last three polynomials define conics:
\[
H_{17}=y(y-\omega x-1)+1,\qquad H_{18}=(y+1+\omega)(y+x-1)+1,\qquad H_{19}=(y+x)(y-(1+\omega)x)+1.
\]

\begin{lemma}\label{lem:conic}
Let $H=(y-ax-b)(y-cx-d)-k$ with $a,b,c,d,k\in\Z[\omega]$, and let $\delta=c-a$. Let $v$ be a finite place at which $\delta$ and $k$ are units, and let $\bar x,\bar y$ be elements of the residue field of $v$ with $\bar H(\bar x,\bar y)=0$. Then $\bar t=\bar y-\bar a\bar x-\bar b$ is non-zero, and
\[
\bar x=\frac{\bar t^2+(\bar b-\bar d)\bar t-\bar k}{\bar\delta\,\bar t},\qquad\bar y=\frac{\bar c\bar t^2+(\bar b\bar c-\bar a\bar d)\bar t-\bar a\bar k}{\bar\delta\,\bar t}.
\]
\end{lemma}

\begin{proof}
We have $\bar t\,(\bar y-\bar c\bar x-\bar d)=\bar k\neq0$, hence $\bar t\neq0$ and $\bar y-\bar c\bar x-\bar d=\bar k/\bar t$. Subtracting this from $\bar y-\bar a\bar x-\bar b=\bar t$ gives $\bar\delta\bar x=\bar t-\bar k/\bar t+\bar b-\bar d$, which is the first formula, and then $\bar y=\bar t+\bar a\bar x+\bar b$ gives the second one.
\end{proof}

For $H_{17}$, $H_{18}$ and $H_{19}$ we have $k=-1$, and $\delta$ is $\omega$, $-1$ and $2+\omega$, respectively. Since $N(2+\omega)=3$, Lemma~\ref{lem:conic} applies at every place not above $3$, and gives the parametrizations
\begin{gather*}
\sigma_{17}(t)=\Bigl(\frac{t^2-t+1}{\omega t},\ t\Bigr),\qquad
\sigma_{18}(t)=\Bigl(-t+2+\omega-\frac1t,\ t-1-\omega\Bigr),\\
\sigma_{19}(t)=\Bigl(\frac{t^2+1}{(2+\omega)t},\ t-\frac{t^2+1}{(2+\omega)t}\Bigr).
\end{gather*}
Put $\Delta_i=1$ for $i<17$, and $\Delta_{17}=\omega t$, $\Delta_{18}=-t$ and $\Delta_{19}=(2+\omega)t$. Then
\[
F_i(t)=\Delta_i^4\,f(\sigma_i(t)),\qquad G_{ij}(t)=\Delta_i^2\,H_j(\sigma_i(t))
\]
are polynomials with coefficients in $\Z[\omega]$; the factor $\Delta_i^2$ is enough because all $H_j$ have total degree at most $2$.

\paragraph{Vertex identities.} For every $i$, let $s_i=\sum_j\rho_{ij}$, where $\rho_{ij}\in\{0,1,2\}$. Then
\begin{equation}\label{eq:vertex}
\Delta_i^{2s_i}\,W_i^3-c_i\prod_{j}G_{ij}^{\rho_{ij}}=F_i\,Q_i,\qquad0\le i\le19,
\end{equation}
where the polynomials $W_i$ are listed in Table~\ref{tab:roots}, and the $Q_i$ are polynomials with coefficients in $\Z[\omega]$, obtained by polynomial division of the left-hand side by $F_i$. For example, $F_0=-t^4-t-2-2\omega$, the neighbours of vertex $0$ are $8$, $13$, $16$ and $17$, with $\rho_{0j}=1,1,1,2$, and $G_{0,8}=t-1$, $G_{0,13}=t-1-\omega$, $G_{0,16}=-1-t-t\omega$ and $G_{0,17}=2+2\omega-t$. With $W_0=2t^3+4$, identity \eqref{eq:vertex} at vertex $0$ reads
\begin{multline*}
(2t^3+4)^3-4(t-1)(t-1-\omega)(-1-t-t\omega)(2+2\omega-t)^2\\
=(-t^4-t-2-2\omega)\bigl(-8t^5-40t^2+12t+(12t+24)\omega\bigr).
\end{multline*}

\begin{table}[htbp]
\centering\footnotesize
\setlength{\tabcolsep}{3pt}
\renewcommand{\arraystretch}{1.16}
\caption{The polynomials $W_i(t)=(P_i(t)+\omega Q_i(t))/m_i$ in the vertex identities \eqref{eq:vertex}.}\label{tab:roots}
\begin{tabular}{@{}rr@{\hspace{9pt}}p{0.425\textwidth}@{\hspace{7pt}}p{0.425\textwidth}@{}}
			\toprule $i$ & $m_i$ & $P_i(t)$ & $Q_i(t)$\\\midrule
		0 & 1 & $2 t^{3} + 4$ & $0$\\\addlinespace[3pt]
		1 & 1 & $- 4 t - 4$ & $- 4 t - 4$\\\addlinespace[3pt]
		2 & 1 & $2 t^{2}$ & $2 t$\\\addlinespace[3pt]
		3 & 1 & $2 t^{2} + 4 t - 8$ & $6 t - 4$\\\addlinespace[3pt]
		4 & 1 & $2 t^{2} - 2 t + 2$ & $2 t^{2} - 2 t + 2$\\\addlinespace[3pt]
		5 & 1 & $2 t - 2$ & $- 2$\\\addlinespace[3pt]
		6 & 1 & $2 t^{2} - 4 t + 8$ & $2 t^{2} - 4 t + 2$\\\addlinespace[3pt]
		7 & 1 & $1$ & $0$\\\addlinespace[3pt]
		8 & 1 & $- 2 t^{3} - 2 t^{2} + 4 t$ & $2 t^{3} - 4 t^{2} + 2 t$\\\addlinespace[3pt]
		9 & 1 & $4 t^{3} - 16 t^{2} + 24 t - 12$ & $2 t^{3} - 8 t^{2} + 12 t - 6$\\\addlinespace[3pt]
		10 & 1 & $2 t^{4} - 2 t - 2$ & $- 2 t^{3} - 4 t^{2} - 2 t - 2$\\\addlinespace[3pt]
		11 & 1 & $- t + 1$ & $- t^{3} + t^{2}$\\\addlinespace[3pt]
		12 & 1 & $4 t^{3} - 8 t^{2} + 8 t - 4$ & $2 t^{3} - 4 t^{2} + 4 t - 2$\\\addlinespace[3pt]
		13 & 1 & $- t^{3} + 1$ & $t^{2} + t + 1$\\\addlinespace[3pt]
		14 & 1 & $2$ & $2$\\\addlinespace[3pt]
		15 & 1 & $3 t^{3} - 2 t^{2} - 2 t - 1$ & $- 4 t^{2} + 2 t - 2$\\\addlinespace[3pt]
		16 & 1 & $- t^{3} - t^{2} - t$ & $- 2 t^{3} - 2 t^{2} - 2 t$\\\addlinespace[3pt]
		17 & 1 & $\begin{aligned}[t]&- t^{7} + t^{6} - 4 t^{5} - t^{4}\\&\quad + 3 t^{3} - 4 t^{2} + 6 t - 4\end{aligned}$ & $\begin{aligned}[t]&- t^{7} + 2 t^{6} - 3 t^{5} - t^{4}\\&\quad + 2 t^{3} - 7 t^{2} + 4 t - 4\end{aligned}$\\\addlinespace[3pt]
		18 & 1 & $\begin{aligned}[t]&- 5 t^{7} + 28 t^{6} - 32 t^{5} - 16 t^{4}\\&\quad + 47 t^{3} - 45 t^{2} + 18 t - 1\end{aligned}$ & $\begin{aligned}[t]&- 4 t^{7} + 41 t^{6} - 125 t^{5} + 176 t^{4}\\&\quad - 164 t^{3} + 98 t^{2} - 40 t + 8\end{aligned}$\\\addlinespace[3pt]
		19 & 3 & $\begin{aligned}[t]&- t^{7} - t^{5} + 9 t^{4} + 18 t^{3}\\&\quad + 8 t^{2} - 3 t + 2\end{aligned}$ & $\begin{aligned}[t]&- 2 t^{7} - 3 t^{6} - 2 t^{5} + t^{2}\\&\quad - 12 t - 2\end{aligned}$\\\addlinespace[3pt]
			\bottomrule
		\end{tabular}
\end{table}

\paragraph{Separation.} For every edge $(i,j)$ with $i<j$, the norm of the resultant $\operatorname{Res}_t(F_i,G_{ij})$ is non-zero and of the form $2^a3^b$, see Table~\ref{tab:resultants}. Also the norms of all $c_i$ are of the form $2^a3^b$.

\begin{table}[htbp]
\centering\small
\setlength{\tabcolsep}{4pt}
\renewcommand{\arraystretch}{1.12}
\caption{Norms of the resultants. The entry $(i,j;a,b)$ means that the norm of $\operatorname{Res}_t(F_i,G_{ij})$ is $2^a3^b$.}\label{tab:resultants}
\begin{tabular}{@{}llll@{}}
			\toprule
			$(i,j;a,b)$ & $(i,j;a,b)$ & $(i,j;a,b)$ & $(i,j;a,b)$\\\midrule
		$(0,8;2,1)$ & $(2,18;0,0)$ & $(5,18;2,0)$ & $(10,11;0,0)$\\
		$(0,13;2,0)$ & $(2,19;0,0)$ & $(5,19;0,0)$ & $(10,13;8,0)$\\
		$(0,16;2,1)$ & $(3,9;4,2)$ & $(6,8;2,2)$ & $(10,17;4,0)$\\
		$(0,17;4,2)$ & $(3,10;2,0)$ & $(6,9;0,1)$ & $(11,16;0,2)$\\
		$(1,2;2,0)$ & $(3,11;0,1)$ & $(6,10;2,0)$ & $(12,15;0,4)$\\
		$(1,4;2,0)$ & $(3,12;2,0)$ & $(6,18;4,0)$ & $(13,15;2,0)$\\
		$(1,5;2,0)$ & $(3,16;4,1)$ & $(6,19;0,0)$ & $(13,18;2,0)$\\
		$(1,6;2,0)$ & $(3,18;6,0)$ & $(7,11;0,0)$ & $(13,19;0,0)$\\
		$(1,10;4,0)$ & $(4,10;2,0)$ & $(7,17;0,0)$ & $(14,15;4,0)$\\
		$(1,12;2,0)$ & $(4,12;2,0)$ & $(8,9;4,4)$ & $(15,17;0,0)$\\
		$(1,13;2,0)$ & $(4,14;4,0)$ & $(8,16;2,4)$ & $(15,18;4,0)$\\
		$(1,15;4,0)$ & $(4,19;0,0)$ & $(8,17;2,1)$ & $(16,19;2,0)$\\
		$(1,17;2,0)$ & $(5,8;2,0)$ & $(9,12;4,0)$ & $(18,19;2,0)$\\
		$(2,9;4,0)$ & $(5,11;0,0)$ & $(9,14;4,0)$ & \\
		$(2,14;4,0)$ & $(5,12;4,0)$ & $(9,15;4,0)$ & \\
			\bottomrule
		\end{tabular}
\end{table}

\begin{proposition}\label{prop:good}
Let $(x,y)$ be an integer solution of \eqref{eq:main}. Then $B_p=0$ for every prime $p\neq2,3$.
\end{proposition}

\begin{proof}
Let $v$ be a place above $p\neq2,3$, and let $k_v$ be its residue field. Let $(\bar x,\bar y)$ be the reduction of $(x,y)$ modulo $v$, and suppose that $v$ divides $H_i$ for some $i$. Then $\bar H_i(\bar x,\bar y)=0$, and there is $\bar t\in k_v$ with $(\bar x,\bar y)=\sigma_i(\bar t)$ and $\bar\Delta_i(\bar t)\neq0$: for $i<17$ we take $\bar t=\bar x$ or $\bar t=\bar y$, and for $i\geq17$ we use Lemma~\ref{lem:conic}. Since $f(x,y)=0$, we get
\begin{equation}\label{eq:reduction}
\bar F_i(\bar t)=0,\qquad\bar G_{ij}(\bar t)=\bar\Delta_i(\bar t)^2\,\bar H_j(\bar x,\bar y)\quad\text{for all }j.
\end{equation}

(a) \emph{If $(i,j)$ is an edge, then $v$ does not divide both $H_i$ and $H_j$.} Indeed, let $i<j$, and suppose that $v$ divides $H_i$ and $H_j$. By \eqref{eq:reduction}, $\bar t$ is a common root of $\bar F_i$ and $\bar G_{ij}$. The resultant can be written as $\operatorname{Res}_t(F_i,G_{ij})=AF_i+BG_{ij}$ with polynomials $A,B$ with coefficients in $\Z[\omega]$, hence $v$ divides $\operatorname{Res}_t(F_i,G_{ij})$. This is impossible, because the norm of this resultant is of the form $2^a3^b$ by Table~\ref{tab:resultants}.

(b) \emph{If $v$ divides $H_i$, then $\bar c_i\prod_j\bar H_j^{\,\rho_{ij}}=\bar W_i(\bar t)^3$ in $k_v$.} Indeed, the coefficients of $Q_i$ are in $\Z[\omega]$, and the coefficients of $W_i$ have denominators dividing $3$. Hence we may reduce \eqref{eq:vertex} modulo $v$ and substitute $\bar t$, and \eqref{eq:reduction} gives
\[
\bar\Delta_i(\bar t)^{2s_i}\Bigl(\bar W_i(\bar t)^3-\bar c_i\prod_j\bar H_j^{\,\rho_{ij}}\Bigr)=0.
\]

Now consider the sum \eqref{eq:B}. All $c_i$ are units at $v$, because their norms are not divisible by $p$. By (H2), a term $e_{ij}\beta_v(H_i,H_j)$ in which both $H_i$ and $H_j$ are units vanishes, and so does a term $\beta_v(H_i,c_i)$ in which $H_i$ is a unit. By (a), every remaining term involves exactly one value $H_i$ divisible by $v$. For such an $i$, we collect all these terms, using (H1) and $\rho_{ij}=-e_{ji}$ for $j<i$:
\[
\sum_{j>i}e_{ij}\beta_v(H_i,H_j)+\sum_{j<i}e_{ji}\beta_v(H_j,H_i)+\beta_v(H_i,c_i)
=\beta_v\Bigl(H_i,\ c_i\prod_jH_j^{\rho_{ij}}\Bigr).
\]
By (a), the second argument is a unit, and by (b) its reduction is a non-zero cube in $k_v$, so that its value of $\chi$ is $0$. By (H2), this symbol is $0$. Hence $B_v=0$ for every place $v$ above $p$, and $B_p=0$.
\end{proof}

\section{\texorpdfstring{The primes $2$ and $3$}{The primes 2 and 3}}\label{sec:local}

Let $p\in\{2,3\}$. We split the hypothetical integer solutions into finitely many residue classes $x\equiv a$, $y\equiv b\pmod{p^k}$, and compute $B_p$ in each class.

\begin{lemma}\label{lem:lifting}
Let $k\geq1$, and let $f(a,b)\equiv0\pmod{p^k}$. Every integer solution of \eqref{eq:main} with $x\equiv a$, $y\equiv b\pmod{p^k}$ satisfies $x\equiv a+p^ks$, $y\equiv b+p^kt\pmod{p^{k+1}}$ for some $0\le s,t<p$ with $f(a+p^ks,b+p^kt)\equiv0\pmod{p^{k+1}}$.
\end{lemma}

\begin{proof}
The digits $s$ and $t$ of $x$ and $y$ in base $p$ are determined by $x$ and $y$, and $f(x,y)=0$ is divisible by $p^{k+1}$.
\end{proof}

\paragraph{Classes of the auxiliary values.} Let $(x,y)$ be an integer solution with $x\equiv a$, $y\equiv b\pmod{p^k}$. Then $H_i(x,y)\equiv H_i(a,b)$ modulo $p^k\Z[\omega]$, and this determines a set of possible classes of $H_i(x,y)$ modulo cubes at the place above $p$, as follows.
\begin{itemize}
\item $p=2$. If $2^k\nmid H_i(a,b)$, then $e=v(H_i(x,y))=v(H_i(a,b))<k$, and the unit part of $H_i(x,y)$ is known modulo $2^{k-e}$, hence modulo $2$, so that the class $(e\bmod3,\chi)$ is determined. If $2^k\mid H_i(a,b)$, we allow all classes, with $\chi=0$ if $H_i$ has rational coefficients.
\item $p=3$. Let $e=v(H_i(a,b))$, the valuation at $\pi$. If $e<2k$, then $v(H_i(x,y))=e$, and the unit part is known modulo $\pi^{2k-e}$. If $2k-e\geq4$, its class is determined by Lemma~\ref{lem:M3}, and otherwise we allow the classes of all units modulo $9$ which are congruent to it modulo $\pi^{2k-e}$. If $e\geq2k$, we allow all classes. If $H_i$ has rational coefficients, we use instead that $H_i(x,y)$ is a rational integer: if $3^k\nmid H_i(a,b)$, its $3$-adic valuation $j$ and its unit part modulo $3^{k-j}$ are known, and we allow the classes of $3^ju$ for all rational units $u$ modulo $9$ compatible with it; if $3^k\mid H_i(a,b)$, we allow the classes of $3^ju$ for all $j\in\{0,1,2\}$ and all units $u$ modulo $9$.
\end{itemize}
In every case, the allowed set contains the class of $H_i(x,y)$. The constants $c_i$ have fixed classes.

\paragraph{The computation.} We start with the solutions of \eqref{eq:main} modulo $p$, and refine a residue class by Lemma~\ref{lem:lifting} as long as \eqref{eq:B}, computed with \eqref{eq:beta2} for $p=2$ and with Lemma~\ref{lem:M3} for $p=3$, takes more than one value when the classes of the $H_i$ run over the allowed sets. Modulo $2$, the solutions are $(x,y)\equiv(0,0)$ and $(1,0)$. The class $(1,0)$ is terminal; the class $(0,0)$ lifts to $(2,0)$ and $(0,2)$ modulo $4$, the first of these is terminal, and the second lifts to $(0,6)$ and $(4,2)$ modulo $8$. Modulo $3$, the only solution is $(1,1)$, which lifts to $(1,7)$, $(4,1)$ and $(7,4)$ modulo $9$, and then to the nine classes modulo $27$ in Table~\ref{tab:local}. By Lemma~\ref{lem:lifting}, every integer solution lies in one of these $13$ terminal classes.

\begin{table}[htbp]
\centering\small
\setlength{\tabcolsep}{5pt}
\caption{The terminal residue classes $x\equiv a$, $y\equiv b\pmod{p^k}$ and the values of $B_p$. The last column lists the $H_i$ whose class is not determined by the residue class. In every class, $B_p$ takes the given value for all allowed choices of these classes.}\label{tab:local}
\begin{tabular}{@{}rrrrrl@{}}
			\toprule
			$p$ & $k$ & $a$ & $b$ & $B_p$ & Undetermined\\\midrule
			2 & 1 & 1 & 0 & 2 & $H_4,H_{12}$\\
			2 & 2 & 2 & 0 & 2 & --\\
			2 & 3 & 0 & 6 & 2 & $H_2$\\
			2 & 3 & 4 & 2 & 2 & --\\
			3 & 3 & 1 & 7 & 0 & --\\
			3 & 3 & 4 & 10 & 0 & $H_{11}$\\
			3 & 3 & 7 & 22 & 0 & $H_{16}$\\
			3 & 3 & 10 & 16 & 0 & --\\
			3 & 3 & 13 & 19 & 0 & $H_{11}$\\
			3 & 3 & 16 & 4 & 0 & $H_{16}$\\
			3 & 3 & 19 & 25 & 0 & --\\
			3 & 3 & 22 & 1 & 0 & $H_{11}$\\
			3 & 3 & 25 & 13 & 0 & $H_{16}$\\
			\bottomrule
		\end{tabular}
\end{table}

In every terminal class, at most two of the $H_i$ have more than one allowed class, and there are at most $9$ choices. For example, in the class $x\equiv1$, $y\equiv0\pmod2$, the classes $(e,\chi)$ of the values at the place above $2$ are
\[
\begin{array}{c|cccccccccc}
i&0&1&2&3&4&5&6&7&8&9\\\hline
H_i&(0,2)&(0,2)&(0,0)&(0,1)&(e_4,0)&(0,2)&(0,2)&(0,1)&(0,2)&(0,0)\\
c_i&(2,0)&(2,2)&(2,1)&(1,1)&(2,0)&(1,1)&(1,0)&(0,0)&(2,0)&(1,0)\\[4pt]
i&10&11&12&13&14&15&16&17&18&19\\\hline
H_i&(0,2)&(0,0)&(e_{12},0)&(0,0)&(0,1)&(0,2)&(0,0)&(0,0)&(0,0)&(0,1)\\
c_i&(1,0)&(0,1)&(2,0)&(0,1)&(1,1)&(1,2)&(1,1)&(0,0)&(0,2)&(0,2)
\end{array}
\]
where $e_4$ and $e_{12}$ are the unknown valuations of $H_4=x+1$ and $H_{12}=x-y+1$ modulo $3$. By \eqref{eq:beta2}, the terms of \eqref{eq:B} which contain $e_4$ add up to $e_4$ times
\[
\sum_j\rho_{4j}\chi(H_j)+\chi(c_4)=2\cdot2+2\cdot2+1\cdot0+2\cdot1+2\cdot1+0=12\equiv0,
\]
where the neighbours of vertex $4$ are $1,10,12,14,19$; similarly the terms which contain $e_{12}$ cancel, and \eqref{eq:B} gives $B_2=2$.

The result of the computation is
\begin{equation}\label{eq:local}
B_2=2,\qquad B_3=0
\end{equation}
at every integer solution of \eqref{eq:main}.


\section{End of the proof}

Suppose that $(x,y)$ is an integer solution of \eqref{eq:main}. By Section~\ref{sec:nonzero}, all values $H_i(x,y)$ are non-zero. By Proposition~\ref{prop:good} and \eqref{eq:local},
\[
\sum_pB_p=B_2+B_3=2\neq0\quad\text{in }\Z/3\Z,
\]
which contradicts Proposition~\ref{prop:reciprocity}. Hence equation \eqref{eq:main} has no integer solutions.
\qed

\begin{thebibliography}{9}
\bibitem{N} J.~Neukirch, \emph{Algebraic Number Theory}, Grundlehren der mathematischen Wissenschaften 322, Springer, Berlin, 1999.
\end{thebibliography}

\end{document}
